\section*{Capítulo \ref{ch:b2:schrodinger-wave-functions} --- A equação de Schrödinger e as funções de onda}

\begin{solution}{exo:b2:schrodinger-wave-functions:1}
$h/\sqrt{2mE}$: \qty{0.123}{nm}, \qty{12.3}{pm}, \qty{3.9}{pm} (a \qty{100}{keV} a
energia cinética é um quinto de $mc^2$; o momento relativístico dá
\qty{3.7}{pm}); nêutron \qty{0.18}{nm}; C$_{60}$ \qty{2.8}{pm}; bola de tênis
$\sim 10^{-34}\,\unit{m}$. Todos, menos a bola, já mostraram interferência.
\end{solution}

\begin{solution}{exo:b2:schrodinger-wave-functions:2}
$A = (\pi a^2)^{-1/4}$; $\langle x\rangle = 0$, $\Delta x = a/\sqrt2$; $0.84$; $0.9995$.
\end{solution}

\begin{solution}{exo:b2:schrodinger-wave-functions:3}
A substituição dá $E\varphi = -(\hbar^2/2m)\varphi'' + V\varphi$. $E/h =
\qty{2.4e14}{Hz}$; \qty{4.8e14}{Hz}, \qty{620}{nm}: a luz vermelha que a transição
emite.
\end{solution}

\begin{solution}{exo:b2:schrodinger-wave-functions:4}
Elétron, \qty{1}{nm}: $\tau = \qty{1.7e-14}{s}$, \qty{60}{\micro m} após \qty{1}{ns},
\qty{60}{km} após \qty{1}{s}; \qty{1}{mm}: $\tau = \qty{17}{s}$, inalterado após \qty{1}{ns},
$+0.2\%$ após \qty{1}{s}; próton, \qty{1}{nm}: $\tau = \qty{3.2e-11}{s}$, \qty{30}{nm},
\qty{30}{m}; grão: $\tau = 600$ anos. Para objetos visíveis $\tau$ é
astronômico --- e o meio ao redor não para de localizá-los.
\end{solution}

\begin{solution}{exo:b2:schrodinger-wave-functions:5}
(a) \cref{prop:b2:schrodinger-wave-functions:current}. (b) $|A|^2\hbar k/m$;
$0$; $(|A|^2 - |B|^2)\hbar k/m$. (c) Os fluxos incidente e refletido
se subtraem sem termo de interferência: os \omterm{thm:b2:wave-interfaces:normal}{coeficientes de transmissão} e de
reflexão podem ser definidos a partir dos fluxos. (d) $\psi^*\psi'$ é real e sua
parte imaginária é nula.
\end{solution}

\begin{solution}{exo:b2:schrodinger-wave-functions:6}
(a) $A = (2\pi\sigma^2)^{-1/4}$, $\langle x\rangle = 0$, $\Delta x = \sigma$. (b) $\hbar k_0$.
(c) $\sigma \times \hbar/2\sigma = \hbar/2$. (d) $\dd\langle x\rangle/\dd t = \int x\,
\partial_t\rho = -\int x\,\partial_xj = \int j = (\hbar/m)\operatorname{Im}\int\psi^*
\psi' = \langle p\rangle/m$.
\end{solution}

\begin{solution}{exo:b2:schrodinger-wave-functions:7}
(a) \qty{5.9e6}{m/s}, \qty{0.123}{nm}, \qty{51}{ns}. (b) $v_\varphi = \qty{3.0e6}{m/s}$,
$v_{\text{g}} = \qty{5.9e6}{m/s}$: a segunda. (c) $\tau = \qty{1.7}{\micro s} \gg
\qty{51}{ns}$: inalterado. (d) $\lambda/a = \qty{4e-7}{rad}$, \qty{0.1}{\micro m}:
os raios servem bem.
\end{solution}

\begin{solution}{exo:b2:schrodinger-wave-functions:8}
(a) $|\psi|^2$ e todos os valores esperados ficam inalterados. (b) $\tfrac12(|\psi_1|^2
+ |\psi_2|^2) + \operatorname{Re}(\eu^{\iu\alpha}\psi_1^*\psi_2)$. (c) As ondas vindas
das duas fendas; sua fase relativa vale $k(r_2 - r_1)$. (d) O detector
termina em estados diferentes para os dois caminhos; as duas componentes já
não interferem, e o padrão é a soma de dois padrões de fenda
única.
\end{solution}

\begin{solution}{exo:b2:schrodinger-wave-functions:9}
(a) $\langle V'\rangle = m\omega^2\langle x\rangle$: exato. (b) $\langle V'\rangle
= mg$: exato. (c) Quando o pacote é estreito frente à escala em que
$V''$ varia. (d) Seu comprimento de onda a \qty{0.1}{m/s} vale $10^{-29}\,\unit{m}$; o
pacote teria de ser tão largo quanto a tigela para importar. O comprimento de onda
térmico de um átomo, \qty{0.1}{nm}, é do tamanho dele próprio.
\end{solution}

\begin{solution}{exo:b2:schrodinger-wave-functions:10}
(a) $\tfrac12(\varphi_1^2 + \varphi_2^2) + \varphi_1\varphi_2\cos\omega_{21}t$. (b) $L/2 +
\cos\omega_{21}t\int x\varphi_1\varphi_2 = L/2 - (16L/9\pi^2)\cos\omega_{21}t$. (c)
$3h^2/8mL^2 = \qty{1.13}{eV}$, \qty{2.7e14}{Hz}, \qty{1.1}{\micro m}. (d) Ele irradia
em $\nu_{21}$: a linha; um \omterm{thm:b2:schrodinger-wave-functions:stationary}{estado estacionário} puro não tem \omterm{thm:b2:dipole-radiation:field}{dipolo
oscilante} e não irradia.
\end{solution}

\begin{solution}{exo:b2:schrodinger-wave-functions:11}
(a) $\psi \propto \int\eu^{-\sigma_0^2k^2}\eu^{\iu(kx - \hbar k^2t/2m)}\dd k$. (b) $\alpha =
\sigma_0^2 + \iu\hbar t/2m$, $\beta = \iu x$: $|\psi|^2 \propto \exp[-x^2/2\sigma_0^2(1 +
(\hbar t/2m\sigma_0^2)^2)]$. (c) $\Delta v = \hbar/2m\sigma_0$, e $\Delta v\,\tau =
\sigma_0$. (d) $\omega = ck$: sem dispersão, todas as componentes a $c$.
\end{solution}

\begin{solution}{exo:b2:schrodinger-wave-functions:12}
(a) $\Gamma = \hbar/\tau$: \qty{6.6e-8}{eV}; \qty{6.6e-8}{eV}; \qty{66}{MeV}. (b)
$1/2\pi\tau = \qty{16}{MHz}$, cem vezes menos que a Doppler. (c) Em
átomos frios e íons aprisionados a largura Doppler desaparece e a largura
natural é atingida --- a base dos padrões de frequência. (d) $c\tau =
\qty{3}{m}$, o comprimento do trem.
\end{solution}

\begin{solution}{pb:b2:schrodinger-wave-functions:1}
\textbf{1.} $v = \sqrt{2eU/m} = \qty{8.4e7}{m/s}$; $p = \qty{7.6e-23}{kg.m/s}$; $\lambda =
\qty{8.7}{pm}$.

\textbf{2.} Erros de alguns por cento: aceitável para estimativas ($\lambda =
\qty{8.6}{pm}$ exatamente).

\textbf{3.} \qty{4.8}{ns}.

\textbf{4.} $\lambda/a = \qty{3e-8}{rad}$: \qty{12}{nm} --- $10^4$ vezes menor que um
pixel.

\textbf{5.} $\Delta\theta = \Delta p_y/p \sim \hbar/2ap = \lambda/4\pi a$: a mesma ordem.

\textbf{6.} $\tau = \qty{1.6}{ms} \gg \qty{4.8}{ns}$: inalterado.

\textbf{7.} $\dd\langle p_y\rangle/\dd t = eE$ exatamente para um campo uniforme:
$\langle y\rangle$ segue a parábola clássica; o tubo é um
instrumento clássico.

\textbf{8.} $v/2 = \qty{4.2e7}{m/s}$; só a \omterm{prop:b2:dispersion-wave-packets:group}{velocidade de grupo} é observável;
a fase de uma onda plana não carrega sinal (e depende do zero
arbitrário de energia).

\textbf{9.} $\Delta p = \hbar/2\sigma_0$, $\Delta v = \hbar/2m\sigma_0$.

\textbf{10.} \qty{1.7e-16}{s}, \qty{1.7e-14}{s}, \qty{1.7e-8}{s}; próton \qty{3.2e-11}{s};
vírus \qty{20}{s}.

\textbf{11.} Ele está num \omterm{thm:b2:schrodinger-wave-functions:stationary}{estado estacionário}: o potencial segura o
pacote; $|\varphi|^2$ não se move.

\textbf{12.} $v = \qty{5.9e5}{m/s}$; a largura dobra em $t = \sqrt3\tau = \qty{2.9e-14}{s}$:
\qty{17}{nm}, uns catorze comprimentos de onda.

\textbf{13.} Suas \omterm{def:b2:schrodinger-wave-functions:psi}{funções de onda} se estendem por todo o cristal: um
momento definido, nenhuma posição definida.

\textbf{14.} Resolução $\sim\lambda/\alpha = \qty{4}{pm}/0.01 = \qty{0.4}{nm}$: são as
aberrações das lentes magnéticas que limitam a abertura, não o comprimento de onda.

\textbf{15.} Compare $\lambda$ com os tamanhos em jogo (aberturas, a
escala de $V$), e o tempo de alargamento com o tempo de observação.

\textbf{16.} \qty{0.38}{}, \qty{1.50}{}, \qty{3.38}{eV}; $\nu_{21} = \qty{2.7e14}{Hz}$
(\qty{1.1}{\micro m}), $\nu_{31} = \qty{7.3e14}{Hz}$ (\qty{410}{nm}).

\textbf{17.} O fasor tem módulo um; em $x = L/2$.

\textbf{18.} Em $t = 0$ a densidade se acumula na metade esquerda; em
$T_{21}/2$, na direita.

\textbf{19.} $\langle x\rangle = L/2 - (16L/9\pi^2)\cos\omega_{21}t$: amplitude
$0.18L = \qty{0.18}{nm}$.

\textbf{20.} Em $\nu_{21}$: a linha espectral; um \omterm{thm:b2:schrodinger-wave-functions:stationary}{estado estacionário} não tem
\omterm{thm:b2:dipole-radiation:field}{dipolo oscilante} e é estável; a superposição irradia até
descer a $\varphi_1$.

\textbf{21.} Não é definida: $\langle E\rangle = (E_1 + E_2)/2 = \qty{0.94}{eV}$; uma
medida dá $E_1$ ou $E_2$, cada um com probabilidade um meio.

\textbf{22.} Aleatório: onde cada elétron cai; não aleatório: o
padrão $|\psi|^2$ que as quedas constroem.

\textbf{23.} $j = |A|^2\hbar k/m$: com $|A|^2 = n$ elétrons por unidade de comprimento,
$nv$ elétrons por segundo.

\textbf{24.} De primeira ordem, para que $\psi(x,0)$ seja uma descrição completa
que fixa o futuro; complexa, para que uma energia definida tenha a
única dependência temporal $\eu^{-\iu Et/\hbar}$ e para que $E = p^2/2m$ apareça.

\textbf{25.} $\psi$ é a amplitude cujo quadrado é a probabilidade; ela
obedece à equação linear de Schrödinger; os \omterm{thm:b2:schrodinger-wave-functions:stationary}{estados estacionários} carregam um
fasor e densidades fixas; um pacote livre se move a $p/m$ e se alarga
em $2m\sigma_0^2/\hbar$; Ehrenfest devolve Newton para o que é grande.
\end{solution}
